\section{Problem Setup}

We consider a fluid-filled elastic shell with equatorial semi-axis $a$, polar
semi-axis $c$, thickness $h$, Young's modulus $E$, Poisson ratio $\nu$, wall
density $\rho_w$, fluid density $\rho_f$, fluid bulk modulus $K_f$, and
internal pressure $P_\mathrm{iap}$.  The canonical parameter set is that used in
the Browntone abdominal-shell studies:
\[
  a=\SI{0.18}{\metre},\quad
  c=\SI{0.12}{\metre},\quad
  h=\SI{0.01}{\metre},\quad
  E=\SI{0.1}{\mega\pascal},
\]
\[
  \nu=\num{0.45},\quad
  \rho_w=\SI{1100}{\kilo\gram\per\metre\cubed},\quad
  \rho_f=\SI{1020}{\kilo\gram\per\metre\cubed},\quad
  K_f=\SI{2.2}{\giga\pascal},\quad
  P_\mathrm{iap}=\SI{1000}{\pascal}.
\]
The corresponding aspect ratio is $\zeta=c/a=\num{0.667}$ and the
equivalent-sphere radius is
$R_\mathrm{eq}=(a^2c)^{1/3}=\SI{0.157}{\metre}$ \cite{BrowntoneP1,BrowntoneP8}.

The coupled shell--fluid model is treated by a Rayleigh--Ritz approximation.
With Ritz coefficients $\mathbf{q}$, the discrete potential and kinetic
energies are
\[
  \Pi=\tfrac{1}{2}\mathbf{q}^\mathsf{T}K(a,c,E)\mathbf{q},
  \qquad
  T=\tfrac{1}{2}\mathbf{q}^\mathsf{T}M(a,c)\mathbf{q},
\]
where $K$ collects the shell stiffness and fluid-compression contributions, and
$M$ collects the shell and added-fluid inertia terms.  Stationarity of
$\Pi-\omega^2T$ gives the generalised eigenproblem
\[
  K(a,c,E)\mathbf{q}=\omega^2 M(a,c)\mathbf{q},
\]
from which the resonance frequencies follow as $f_n=\omega_n/(2\pi)$.
Papers~1 and 8 give the full derivations of $K$ and $M$
\cite{BrowntoneP1,BrowntoneP8}, but these equations define the forward model
used here.

The inverse variable is the parameter vector
$\boldsymbol{\theta}=[a,c,E]^\mathsf{T}$.  The spectral data vector is formed
from the first five flexural resonances,
\[
  \mathbf{f}(\boldsymbol{\theta})
  =
  [f_2,f_3,f_4,f_5,f_6]^\mathsf{T},
\]
where $f_n$ is the resonance frequency of flexural mode $n$.  At the canonical
oblate point, the lowest member of this set is
$f_2=\SI{3.95}{\hertz}$ \cite{BrowntoneP1}.  We do not use the breathing mode in
this inverse map.

Local practical identifiability is measured through the dimensionless scaled
Jacobian
\[
  \widetilde{J}
  =
  \operatorname{diag}(\mathbf{f})^{-1}
  J
  \operatorname{diag}(\boldsymbol{\theta}),
  \qquad
  \widetilde{J}_{ij}
  =
  \frac{\theta_j}{f_i}
  \frac{\partial f_i}{\partial \theta_j},
  \qquad
  J_{ij}=\frac{\partial f_i}{\partial \theta_j},
\]
with singular values $\sigma_1\geq\sigma_2\geq\sigma_3\geq 0$.  The
conditioning metric is
\[
  \kappa=\frac{\sigma_1}{\sigma_3},
\]
provided $\sigma_3>0$.  Large $\kappa$ indicates a weakly observable parameter
combination; $\sigma_3=0$ implies local rank deficiency.

Three geometry classes are compared.
\begin{enumerate}
  \item \textbf{Equivalent sphere.}  The forward model depends on $a$ and $c$
  only through $R_\mathrm{eq}$.  The geometry columns of the exact Jacobian are
  therefore linearly dependent, so the map is singular in the strict algebraic
  sense \cite{BrowntoneP7,BrowntoneP8}.
  \item \textbf{Oblate Ritz family.}  The full spheroidal curvature field is
  retained.  For the oblate branch we use the eccentricity
  $\varepsilon=\sqrt{1-c^2/a^2}$ and the aspect-ratio offset
  $\delta=1-c/a$, and examine the near-spherical limit together with the
  canonical oblate point.
  \item \textbf{Prolate Ritz family.}  The same five-mode diagnostic is applied
  to a prolate continuation away from the sphere, so that the effect of
  geometry class can be separated from the mere fact of leaving
  $\mathrm{SO}(3)$ symmetry behind.  Along this branch $c>a$, the enclosed
  volume and wall thickness are fixed, and the path is parameterised by
  $c/a\in[\num{1.0},\num{1.5}]$ with $a$ decreasing and $c$ increasing subject
  to $a^2c=\text{const}=R_\mathrm{eq}^3$.
\end{enumerate}
